
%  PLANTILLA GUÍAS INSTITUCIONALES — USACH PREMIUM
%  Fusiona el estilo formal de guia1 (portada + marca de agua)
%  con el estilo premium de guia3 (colores + tcolorbox + multicol).
%
%  INSTRUCCIONES RÁPIDAS:
%  1. Edita el bloque "INFORMACIÓN DE LA GUÍA" (líneas ~65-85).
%  2. Rellena \listacontenidos con los temas de tu guía.
%  3. Agrega ejercicios en los entornos indicados.
%  4. Compila con: pdflatex guia-institucional.tex (2 veces).
%
%  ESTRUCTURA DEL DOCUMENTO:
%    Portada → Introducción → Kit de Propiedades → Ejercicios → Solucionario
% ============================================================

\documentclass[letterpaper, 12pt]{article}

% ============================================================
% PAQUETES
% ============================================================
\usepackage[utf8]{inputenc}
\usepackage[spanish]{babel}
\usepackage{amsmath, amssymb}
\usepackage{geometry}
\usepackage{fancyhdr}
\usepackage{enumitem}
\usepackage{graphicx}
\usepackage{hyperref}
\usepackage{multicol}
\usepackage{parskip}
\usepackage{xcolor}
\usepackage{tcolorbox}
\usepackage{eso-pic}
\usepackage{tikz}
\usetikzlibrary{patterns}
\tcbuselibrary{skins}

% ============================================================
% GEOMETRÍA
% ============================================================
\geometry{letterpaper, margin=1.5cm, top=3cm, bottom=2.5cm, headheight=2cm}

% ============================================================
% COLORES INSTITUCIONALES
% ============================================================
\definecolor{celesteSuave}{HTML}{F0F7FF}
\definecolor{azulFuerte}{HTML}{1A5276}
\definecolor{turquesaUsach}{HTML}{33CCCC}
\definecolor{enlace}{HTML}{1A5276}

% ============================================================
% RUTAS DE IMÁGENES  ← ajusta si tu carpeta de imágenes es distinta
% ============================================================
\newcommand{\logoup}{../../../../../template/images/logo_up.png}
\newcommand{\logousach}{../../../../../template/images/logo_usach.png}
\newcommand{\logofondo}{../../../../../template/images/logo_up.png}

% ============================================================
% INFORMACIÓN DE LA GUÍA  ← EDITA AQUÍ ANTES DE COMPILAR
% ============================================================
\newcommand{\institucion}{Usach Premium}
\newcommand{\curso}{Calculo I}              % ej: Cálculo I
\newcommand{\guiasnumero}{Guía N - Título}         % ej: Guía 1 - Trigonometría
\newcommand{\semestre}{1er. Semestre de 2026}
\newcommand{\preparadopor}{Usach Premium.}         % o el nombre del ayudante
\newcommand{\webinstitucion}{www.usachpremium.cl}
\newcommand{\instagraminstitucion}{https://www.instagram.com/usach.premium}
\newcommand{\transparencia}{0.05}                  % opacidad marca de agua (0.0–1.0)
\newcommand{\fraseportada}{"Por y para estudiantes"}

% Lista de contenidos que aparece en la portada dentro del tcolorbox.
% Agrega o quita \item según los temas de tu guía.
\newcommand{\listacontenidos}{
    \item[\color{azulFuerte}$\Rightarrow$] Inecuaciones con raíz
}

% ============================================================
% HYPERREF
% ============================================================
\hypersetup{
    colorlinks=true,
    linkcolor=enlace,
    urlcolor=enlace,
    citecolor=enlace,
    pdfborder={0 0 0},
}

% ============================================================
% ENCABEZADO Y PIE DE PÁGINA (páginas interiores)
% ============================================================
\pagestyle{fancy}
\fancyhf{}
\fancyhead[L]{%
    \includegraphics[height=1.2cm]{\logoup}%
}
\fancyhead[R]{%
    \begin{minipage}[b]{0.42\textwidth}
        \raggedleft
        \fontsize{9pt}{11pt}\selectfont
        \textbf{\color{azulFuerte}\institucion} \\
        \curso\ — \guiasnumero \\
        \textit{\color{gray}@usach.premium}
    \end{minipage}\hspace{0.1cm}%
    \includegraphics[height=1.2cm]{\logousach}%
}
\fancyfoot[L]{\fontsize{9pt}{11pt}\selectfont \textit{\fraseportada}}
\fancyfoot[R]{\fontsize{9pt}{11pt}\selectfont Página \thepage}
\renewcommand{\headrulewidth}{0.4pt}
\renewcommand{\footrulewidth}{0.4pt}

% Estilo portada: sin encabezado ni pie
\fancypagestyle{plain}{
    \fancyhf{}
    \renewcommand{\headrulewidth}{0pt}
    \renewcommand{\footrulewidth}{0pt}
}

% ============================================================
% MARCA DE AGUA (logo transparente al centro de cada página)
% Llama a \MarcaAgua dentro de un entorno para activarla.
% ============================================================
\newcommand{\MarcaAgua}{
    \AddToShipoutPicture{
        \AtPageCenter{
            \begin{tikzpicture}[remember picture, overlay]
                \node[opacity=\transparencia] at (0,0)
                    {\includegraphics[width=0.7\textwidth]{\logofondo}};
            \end{tikzpicture}%
        }
    }
}

% ============================================================
% CAJAS REUTILIZABLES
% ============================================================

% Caja de resumen / kit de propiedades
% Uso: \begin{kitpropiedades}{Título de la caja} ... \end{kitpropiedades}
\newenvironment{kitpropiedades}[1]{%
    \begin{tcolorbox}[colback=celesteSuave, colframe=azulFuerte,
        title=\textbf{#1}, fonttitle=\bfseries]
}{%
    \end{tcolorbox}
}

% Caja de desafío con sombra y título de color blanco sobre azul
% Uso: \begin{desafiobox}{Título} ... \end{desafiobox}
\newtcolorbox{desafiobox}[1]{
    colback=white,
    colframe=azulFuerte,
    fonttitle=\bfseries,
    coltitle=white,
    title=#1,
    arc=8pt,
    boxrule=1.2pt,
    enhanced,
    drop shadow={black!5}
}

% Caja de bloque de soluciones
% Uso: \begin{solbox}{Título} ... \end{solbox}
\newenvironment{solbox}[1]{%
    \begin{tcolorbox}[colback=celesteSuave, colframe=azulFuerte,
        title=\textbf{#1}, fonttitle=\bfseries]
}{%
    \end{tcolorbox}
}

% ============================================================
% ENTORNOS DE SECCIÓN
% Cada entorno activa la marca de agua en las páginas que lo contienen.
% ============================================================

\newenvironment{introduccion}{%
    \section*{Introducción}%
    \MarcaAgua
}{}

\newenvironment{ejercicios}{%
    \MarcaAgua
}{}

\newenvironment{soluciones}{%
    \newpage
    \begin{center}
        \begin{tcolorbox}[colback=azulFuerte, colframe=azulFuerte,
            colupper=white, width=0.8\textwidth, center, arc=10pt]
            \centering\Large\textbf{SOLUCIONARIO}
        \end{tcolorbox}
    \end{center}
    \MarcaAgua
    \small
}{}

% ============================================================
\begin{document}
\thispagestyle{plain}

% ============================================================
% PORTADA
% ============================================================

% Logos en la portada: UP arriba a la izquierda, USACH arriba a la derecha
\AddToShipoutPicture*{%
    \put(54, 700){\includegraphics[width=2cm]{\logoup}}
    \put(420, 706){\includegraphics[width=5cm]{\logousach}}
}

\vspace*{0.5cm}

\begin{center}
    \color{azulFuerte}
    {\huge \textbf{GUÍA DE EJERCICIOS}} \\[0.5cm]
    {\Huge \textbf{\curso}} \\[0.3cm]
    {\Large \guiasnumero} \\[1.5cm]

    % Caja de contenidos
    \begin{tcolorbox}[colback=celesteSuave, colframe=azulFuerte, arc=10pt,
        width=0.85\textwidth, center, boxrule=1.5pt]
        \centering
        \vspace{0.2cm}
        {\Large \textbf{Contenidos}}
        \vspace{0.3cm}
        \color{black}
        \begin{itemize}[leftmargin=3cm]
            \listacontenidos
        \end{itemize}
        \vspace{0.2cm}
    \end{tcolorbox}
\end{center}

\vspace{1.5cm}

\begin{center}
    {\Large \textbf{\semestre}} \\[0.8cm]
    {\large \textbf{\preparadopor}}
\end{center}

\vspace{2cm}

% Pie de portada con links
\begin{center}
    {\color{azulFuerte}\hrule height 0.5pt}
    \vspace{0.3cm}
    \begin{minipage}{0.45\textwidth}
        \centering
        \textbf{Instagram:} \\
        \small\href{\instagraminstitucion}{@usach.premium}
    \end{minipage}
    \begin{minipage}{0.45\textwidth}
        \centering
        \textbf{Web:} \\
        \small\href{http://\webinstitucion}{\webinstitucion}
    \end{minipage}
\end{center}

\pagenumbering{arabic}
\newpage

% ============================================================
% ============================================================
% EJERCICIOS
% ============================================================
\begin{ejercicios}

% --- Sección con ejercicios cortos (4 columnas) ---
\subsection*{I. Ejercicios Propuestos}
1.1) \textbf{Resolver la inecuación:$$\sqrt{x^2 - 9} \le x - 1$$}

\medskip
1.2) \textbf{Hallar el conjunto solucion}
$$\frac{1}{\sqrt{x - 2}} < 1$$

\medskip

1.3) \textbf{Resuelva la siguiente inecuacion}
 $$x + 1 < \sqrt{x^2 + 3x}$$
 \medskip

1.4) 
\textbf{Determine el conjunto de solucion}
$$\sqrt{x - 1} < \frac{2}{\sqrt{x - 1}}$$
\medskip

1.5)\textbf{Resolver $$\sqrt{x + 2} - \sqrt{x - 1} < 1$$}
\medskip

1.6)\textbf{Encuentre el Conjunto solucion de la siguiente inecuacion}
$$\frac{\sqrt{x^2 - 1}}{x - 2} > 0$$
\vspace{4 cm}
\newpage

% --- Sección con ejercicios medianos (2 columnas) ---
\subsection*{II. Ejercicios Propuestos}
2.1) \textbf{Considere las siguiente inecuacion}
$$\frac{x - 5}{\sqrt{x - 2}} < 0$$




2.2) \textbf{Resuelva}
$$\frac{x - 8}{\sqrt{x - 3}} \leq 0$$
    \medskip
2.3) \textbf{Resuelva la siguiente Inecuacion}
$$\sqrt{x - \sqrt{2-x}} < 1$$
\medskip
2.4) \textbf{Resuelva la siguiente Inecuacion}
$$\frac{1}{\sqrt{3 - \sqrt{x+1}}} > 1$$
\medskip
2.5 \textbf{Resuelva la siguiente Inecuacion}
$$\frac{\sqrt{x - \sqrt{x}}}{x - 3} \geq 0$$
\medskip
\vspace{4 cm}
\newpage

% ============================================================
% SECCION III — RAIZ DENTRO DE RAIZ
% ============================================================
\subsection*{III. Raíz dentro de Raíz}

3.1) \textbf{Resuelva la siguiente inecuación}
$$\sqrt{\sqrt{x} - 1} \leq 2$$
\medskip

3.2) \textbf{Resuelva la siguiente inecuación}
$$\sqrt{2 + \sqrt{x - 1}} < 2$$
\medskip

3.3) \textbf{Encuentre el conjunto solución}
$$\sqrt{x + \sqrt{x - 1}} \geq 2$$
\medskip

3.4) \textbf{Determine el conjunto solución}
$$\sqrt{3 - \sqrt{2 + x}} \geq 1$$
\medskip

3.5) \textbf{Resuelva la siguiente inecuación}
$$\sqrt{1 - \sqrt{1 - x^2}} \leq \frac{\sqrt{2}}{2}$$
\vspace{3cm}
\newpage

\end{ejercicios}

% ============================================================
% SOLUCIONARIO
% ============================================================
\begin{soluciones}

\begin{solbox}{I. Soluciones — Raíz cuadrada}
\begin{multicols}{2}
\begin{enumerate}[label=\color{azulFuerte}\bfseries 1.\arabic*), leftmargin=*, itemsep=0.4em]
    \item $x \in [3,\,5]$
    \item $x \in (3,\,+\infty)$
    \item $x \in (-\infty,\,-3\,] \cup (1,\,+\infty)$
    \item $x \in (1,\,3)$
    \item $x \in (2,\,+\infty)$
    \item $x \in (2,\,+\infty)$
\end{enumerate}
\end{multicols}
\end{solbox}

\begin{solbox}{II. Soluciones — Raíz cuadrada (continuación)}
\begin{multicols}{2}
\begin{enumerate}[label=\color{azulFuerte}\bfseries 2.\arabic*), leftmargin=*, itemsep=0.4em]
    \item $x \in (2,\,5)$
    \item $x \in (3,\,8\,]$
    \item $x \in \left[1,\,\dfrac{1+\sqrt{5}}{2}\right)$
    \item $x \in (3,\,8)$
    \item $x \in \{0,\,1\} \cup (3,\,+\infty)$
\end{enumerate}
\end{multicols}
\end{solbox}

\begin{solbox}{III. Soluciones — Raíz dentro de raíz}
\begin{multicols}{2}
\begin{enumerate}[label=\color{azulFuerte}\bfseries 3.\arabic*), leftmargin=*, itemsep=0.5em]
    \item $x \in [1,\,25]$
    \item $x \in [1,\,5)$
    \item $x \in \left[\dfrac{9 - \sqrt{13}}{2},\,+\infty\right)$
    \item $x \in [-2,\,2]$
    \item $x \in \left[-\dfrac{\sqrt{3}}{2},\,\dfrac{\sqrt{3}}{2}\right]$
\end{enumerate}
\end{multicols}
\end{solbox}

\vspace{0.8cm}
\begin{center}
    \small\textbf{Usach Premium} — ¡Nos vemos en la ayudantía! \\
    \textit{"Por y para estudiantes."}
\end{center}

\end{soluciones}

\end{document}